\section{Related work}
\label{sec:related}

Published savings of wind propulsion spread widely. A systematic review of 48 studies finds a mean CO$_2$ reduction of 17\,\%, with half the studies between 7.5 and 22.5\,\%, and names inconsistent methods, test conditions and reporting of the ship as the cause \citep{huang2025}. Flettner rotors are the most studied and most installed rig \citep{huang2025}. Few studies compare technologies on the same ship, and fewer on cost. \Cref{app:published} gives the figures of the studies below.

\citet{traut2014} linked a towing kite and a Flettner rotor to a year of wind on five trade routes at a fixed speed and found that which was better depended on the route. They left out side force, routing and cost. \citet{lu2020} compared a rotor, a wing sail and a DynaRig, a square-rigged sail on rotating masts, on a tanker at a fixed arrival time, and the rotor saved the most with the least sail area. \citet{rechevilanova2021} predicted the performance of the same three rigs in all six degrees of freedom of the ship's motion and agreed fairly with a year of trial data from the rotor tanker Maersk Pelican. A study for the European Commission found kites relatively better on small ships, with one \SI{400}{m^2} kite fitted to every ship, and calls its own estimates over-optimistic for ignoring side force \citep{cedelft2016}. EMSA computes the saving a rig needs to break even per ship segment and states that kites perform better at higher ship speed \citep{emsa2023}. The voyage simulation of \citet{podeur2016} found the opposite for an \SI{800}{m^2} kite on a container ship, with drift neglected (\cref{tab:published}).

Rotors have the most detailed models. The ShipCLEAN model balances surge, sway, roll and yaw, the forces along and across the ship and the moments about its long and vertical axes, with the hull's lift and drag at leeway \citep{tillig2020}. It predicts fuel savings of 30\,\% for a tanker and 14\,\% for a ro-ro ship, which carries rolling cargo driven on board. Coupled to transport economics, it shows a rotor on the North Atlantic breaking even after 16.7 years at 500 USD/t of fuel \citep{tillig2020b}. \citet{hansen2025} priced four rotors on a Capesize bulk carrier per round trip without routing and found that the mean wind speed decides the case. On a wind-powered cargo ship the resistance the sails induce, mostly from the rudder, is comparable to the added resistance in waves, while moving the centre of effort of the sails changes the resistance by only 1 to 2\,\% \citep{kramer2022}.

Towing kites have their own models. \citet{erhard2013} controlled the SkySails kites with a quasi-steady model that neglects the kite's mass. \citet{penloup2021} varied the kite's lift, glide ratio (lift over drag), tether length and allowed wing load, the force per unit of kite area, and found that the wing load explains more than half of the gain. These models take the kite's lift and drag as given, at glide ratios of 3.5 \citep{traut2014} and 4.6 \citep{podeur2016}. SkySails is the only maker whose ship kites have been in regular service, from 2011 to 2013 \citep{skysails2018}. The record of the EU WINTECC project reports 5\,\% fuel savings on an average route mix on the Beluga SkySails and 10 to 12\,\% on North Atlantic and North Pacific routes \citep{wintecc2009_record}. Rigid kites, stiff tethered wings, have been designed and flown to generate power \citep{echeverri2020part1,bauer2018} and scaled for that use \citep{eijkelhof2020reference,sommerfeld2022scaling,vanhagen2023,joshi2025}. None of these studies puts one on a ship.

Weather routing raises the value of wind. On a multipurpose ship with four rotors it raised the eastbound energy saving from 36 to 53\,\% and the westbound one from 14 to 28\,\% \citep{bentin2016}, and an industry study cited by EMSA raised the saving of three rotors from 9 to 25\,\% \citep{emsa2023}. \citet{mason2023} routed a Panamax bulk carrier with rotors and showed that forecast errors remove part of the routing gain, which re-planning on new forecasts largely restores. \citet{wendling2025} showed that their router converges as grid and weather data are refined and that coarse weather data understate the gain. \citet{dhainaut2025} optimised route, speed and rotor position together on one container ship and name the comparison of wind technologies, heel, waves and profitability as next steps.

Three gaps remain. No study puts kites, wing sails and rotors on one ship with one balance of side force, one router and one cost model. The nearest, \citet{bentin2018}, routed a kite, a rotor and a DynaRig on one multipurpose ship at a fixed \SI{13}{kn}, with no wing sail and no balance of side force by leeway, heel and rudder, and gives only a rough break-even fuel price taken from another study. \citet{hansen2025} count a full voyage cost, but for rotors alone. Next, the kite enters these comparisons as a given force law and size. None prices the wear of its canopy by the hour flown, a running cost the deck rigs do not have, against the deck rigs' capital in a full voyage cost, and EMSA only charges kites a higher maintenance fraction \citep{emsa2023}. Last, no study sizes every rig for every ship, route and fuel, so none can say which rig comes out ahead in which situation. \Cref{sec:res_speed} compares what the studies find on speed with our results.
